problem:  
Let \(G\) be a connected, simply connected **solvable, non‑nilpotent Lie group** with Lie algebra  
\[
\mathfrak g = \mathbb R X_0 \ltimes_{\phi} \mathfrak n,
\]
where \(\mathfrak n\) is the nilradical and \(\phi=\operatorname{ad}(X_0)|_{\mathfrak n}\) is **invertible**, with all eigenvalues having **nonzero real parts**. Thus \(G\) is non‑unimodular and the adjoint action is exponentially expanding or contracting.  

Each left‑invariant Lorentzian metric \(g_B\) on \(G\) corresponds to a symmetric bilinear form \(B\) of signature \((1,\dim G-1)\) on \(\mathfrak g\). In body coordinates, the geodesic equation becomes  
\[
\dot u = -\,\operatorname{ad}(u)^{\top_B}u,\qquad u(t)\in\mathfrak g.
\]
The metric is **geodesically complete** when all solutions \(u(t)\) exist globally for \(t\in\mathbb R\).

**Problem:**  
Determine whether there exist Lorentzian forms \(B\) on such \(\mathfrak g\) satisfying:

1. \(B(X_0,X_0)=0\) (the distinguished generator \(X_0\) is **null**); and  
2. the corresponding metric \(g_B\) is **geodesically complete**.  

In other words:  
Can a non‑unimodular solvable group with an **expanding derivation** admit a **geodesically complete** left‑invariant Lorentzian metric whose null direction is generated by an expanding element \(X_0\)?  

Refined questions:  
- Does completeness depend on whether all real parts of eigenvalues of \(\phi\) are positive (“uniformly expanding”) or of mixed signs (“contracting/expanding combined”)?  
- Can one describe an algebraic obstruction or balance mechanism preventing or enabling global completeness in the Euler‑Arnold ODE, despite exponential adjoint growth?

**Context:**  
Known completeness results for left‑invariant pseudo‑Riemannian metrics fall into settings that do *not* include this one:  
- **Riemannian case:** all matrices \(\operatorname{ad}(u)\) are skew‑adjoint; completeness always holds.  
- **Nilpotent or unimodular pseudo‑Riemannian groups:** Guediri (1998) and Dotti–Freitas (2002) show completeness whenever \(\operatorname{ad}(u)\) is nilpotent or trace‑free.  
- **Bi‑invariant metrics:** geodesically complete by construction.  
Our solvable \(G\) is non‑unimodular and has \(\operatorname{tr}\phi\neq0\), so no known theorem applies. The indefinite Lorentzian signature and the *null generator* assumption introduce degenerate energy directions complicating the usual blow‑up mechanisms in polynomial ODEs.

Why is it a "good" problem:  
This problem directly addresses an **unsettled boundary** between algebraic growth in solvable groups and global Lorentzian geometry. It asks whether the **indefiniteness of the metric** can offset exponential instability arising from the spectrum of \(\phi\)—a source of incompleteness in most positive‑definite settings.  

Its resolution would demand a synthesis of:  
- **Lie‑algebra spectral analysis** (understanding how \(\Re(\lambda(\phi))\) affects flow behavior);  
- **Lorentzian geometric structure** (role of null directions in completeness); and  
- **qualitative ODE theory** for non‑homogeneous quadratic vector fields.  

A positive result would reveal new balance phenomena where indefinite geometry neutralizes algebraic expansion; a negative result would yield a **spectral obstruction theorem**, showing that real eigenvalues of \(\phi\) inevitably cause incompleteness regardless of signature.  

The question is compact—“Can algebraic expansion coexist with Lorentzian completeness in solvable groups?”—yet it reaches across differential geometry, Lie theory, and dynamical systems.  It exemplifies a **simple statement with deep cross‑disciplinary consequences**, a hallmark of a good research problem.